\section{What a small determinant must accomplish}
For $s>1$, the Dirichlet beta function is
\[
 \beta(s)=\sum_{n=0}^{\infty}\frac{(-1)^n}{(2n+1)^s}.
\]
Its values at positive odd integers are nonzero rational multiples of
powers of $\pi$. For even arguments, explicit irrationality results
remain much harder to obtain. Our earlier work proves that at least one
of $\beta(4),\beta(6),\ldots,\beta(32)$ is irrational
\cite{fifteen}, refining the construction of Krisman \cite{Krisman}.
The investigations here sought a smaller set, or the individual value
$\beta(4)$.

The recurring difficulty is easy to state. A rational-coefficient
polynomial can be very small at the target constants because all of
its coefficients share a very small rational factor. That factor has
no bearing on irrationality. What matters is the value after the
coefficients have been made relatively prime integers. Positivity can
prove that this value is nonzero, but cannot make its arithmetic cost
disappear.

\subsection{The exact denominator result}
Our main result concerns a positive moment functional introduced by
Carreira \cite[Section~11]{Carreira}. For a positive integer $K$,
interlacing $K$ integer and $K$ half-integer poles gives matrices
whose entries are rational combinations
of $1$, $\beta(4)$, and $\zeta(4)$. One possible strategy is to retain
a subspace on which the $\zeta(4)$ coefficient is positive, then replace
$\zeta(4)$ by an upper rational approximation. Positivity supplies
nonvanishing automatically. The question becomes whether the resulting
determinant remains small after clearing its denominators.

Sections~\ref{sec:moments}--\ref{sec:substitution} answer that question
for a precise certificate. Theorem~\ref{thm:valuation} determines the
exact dyadic valuation of every principal minor of the beta-coefficient
matrix. On the positive subspace under consideration it forces a
denominator of size at least
\[
 \exp\bigl((\log 2)K^2-O(K\log K)\bigr).
\]
This is a necessary cost of the whole polynomial, rather than the cost
of a convenient way of clearing its entries. Theorem~\ref{thm:proxy}
shows that the usual error bound $q^{-2}$ for a rational approximation
cannot pay that cost in the specified certificate, for any choice of
its denominator $q$. The distinction between this certificate and the
actual approximants is essential and is made explicit there.

The arithmetic mechanism may be useful separately from this application.
A finite binomial average replaces a residue sum with substantial
cancellation. In its Krawtchouk basis, multiplication has a triangular
reduction modulo two. Known orthogonal norms then account for the
entire determinant valuation. All properties of that basis used here
are derived from its generating function.

\subsection{Other constructions and what they establish}
There is another way to remove $\zeta(4)$: restrict the coefficient
quadratic form to a rational subspace on which it vanishes.
Section~\ref{sec:compression} constructs such subspaces with controlled
height. A local factorial cancellation reduces an apparently exponential
vector-height cost to a polynomial one. Meyer and Cassels provide the
classical existence and small-zero inputs; the local cancellation and
integral lift supply the quantitative adaptation.

Section~\ref{sec:content} develops arithmetic savings for the version
that keeps only the added low-degree directions. It combines pole
collisions, exact low-rank terms, and dyadic Cauchy-minor divisibility.
The resulting upper bound for primitive arithmetic remains insufficient
for the available decay comparison. An insufficient upper bound does
not prove that the actual primitive values fail.

Section~\ref{sec:gram} treats a different positive Gram system involving
$\beta(4),\ldots,\beta(28)$. Removing a rational block leaves a
seven-by-seven determinant. We prove that its logarithm has linear,
rather than quadratic, scale in the family index. Its primitive
arithmetic remains unresolved. Section~\ref{sec:linear} returns to
linear forms: formal rank and integral removal of zero multipliers
lead to eleven controlled independent directions, but do not establish
nonvanishing at the actual constants.

These are later investigations than those collected in
\cite{firstsupp,furthersupp}. Positive determinant criteria and their
arithmetic are established methods \cite{Zudilin}; our contribution
lies in the particular reductions, constructions, and estimates proved
here. The separate constructions are presented as case studies, each
with its conclusion and missing input. A reader principally interested
in the exact valuation and its application can read
Sections~\ref{sec:moments}--\ref{sec:substitution} independently.
The prerequisites are elementary $p$-adic valuations, Gram and
Schur-complement linear algebra, and lattice reduction. The finite
orthogonal-polynomial identities needed for the main theorem are
derived below.

\subsection{Primitive normalization and growing degree}
For a nonzero rational polynomial $D$, let $q_D>0$ be the unique rational
number such that $q_DD$ has integer coefficients with gcd one. We retain
the sign of $D$; whenever $D$ is a positive Gram determinant, its primitive
value is positive. A nonzero rational factor in $D$ has no separate
effect on this primitive polynomial, apart from sign.

If $D$ has degree at most $h$ and a target number equals $a/b$, then
$b^h q_DD(a/b)$ is an integer. For fixed $h$, a positive primitive value
tending to zero gives an irrationality contradiction. When $h=O(K)$,
one must also beat every fixed factor $b^h$; a bound
$q_DD(\bet)\leq\exp(-\varepsilon K^2)$ suffices. We will keep this
degree cost visible in each comparison.

For a univariate determinant with leading nonzero coefficient $\ell_D$,
it is often useful to separate the invariant quantities
\begin{equation}\label{eq:invariant}
 A_D=q_D|\ell_D|,\qquad
 \delta_D=\frac{D(\bet)}{|\ell_D|},\qquad
 q_DD(\bet)=A_D\delta_D.
\end{equation}
Here $A_D$ is the absolute leading coefficient of the primitive polynomial.
Both quantities are unchanged by multiplying $D$ by a positive rational
number. This prevents an arbitrary matrix scale from being counted as
either arithmetic improvement or analytic decay.
